
\documentclass[a4paper,12pt]{article}
\usepackage{amsmath}
\usepackage{amssymb}

\begin{document}

\section*{Electric Conduction}

\begin{enumerate}
  \item By how much does the emitted current change if the tungsten cathode temperature is increased from $T = 2200\,\text{K}$ to $2350\,\text{K}$?  
  The work function for tungsten is $W_{\text{ki}} = 4.5\,\text{eV}$.  
  The Richardson coefficient is $6.02 \cdot 10^{5}\,\text{A}/\text{m}^2\text{K}^2$.

  \item What retarding voltage reduces the emission current density to one fifteenth for a cathode at $T = 2200\,\text{K}$?  
  The electron charge is $q = -1.6 \cdot 10^{-19}\,\text{As}$.

  \item What is the free-electron drift velocity in a copper conductor of radius $R = 1.5\,\text{mm}$ if the current is $I = 12\,\text{A}$?  
  The electron density in copper is $n_{\text{e}} = 3.3 \cdot 10^{28}\,\text{m}^{-3}$.

  \item We examine an n-type semiconductor under the conditions shown in the figure:  
  $B_{y} = 0.4\,\text{T}$, $U_{\text{H}} = 10\,\text{mV}$, $I = 30\,\text{mA}$.  
  The magnetic induction is directed along the $y$–axis.

  \begin{enumerate}
    \item What is the electron concentration?
    \item What is the drift velocity of the electrons?
    \item What is the Hall constant?
  \end{enumerate}

\end{enumerate}

\end{document}
